\begin{pastexam}[label={px:ncku-107-cce-09}]{成大電通所 107 學年度 問答題第 5 題}
（問答；14 分）

\textbf{5. (14 pts) Spanning Tree:}

Find the minimal cost spanning tree according to (a) Kruskal's algorithm and (b) the Sollin's Algorithm, respectively. Please show the result in each iteration of the algorithm.

\begin{center}
\begin{tikzpicture}[v/.style={circle,draw,minimum size=6mm,inner sep=0pt,font=\small\sffamily},w/.style={font=\small\sffamily,inner sep=1.5pt},>={Stealth[length=2mm]}]
\node[v] (0) at (9.38,-19.71) {0};
\node[v] (5) at (7.47,-20.82) {5};
\node[v] (1) at (11.34,-20.77) {1};
\node[v] (6) at (9.43,-21.78) {6};
\node[v] (4) at (7.52,-22.84) {4};
\node[v] (2) at (11.43,-22.65) {2};
\node[v] (3) at (9.50,-23.69) {3};
\draw[] (0) to node[w,above left] {58} (5);
\draw[] (0) to node[w,above right] {20} (1);
\draw[] (5) to node[w,above right] {10} (6);
\draw[] (1) to node[w,above left] {32} (6);
\draw[] (1) to node[w,right] {44} (2);
\draw[] (5) to node[w,left] {25} (4);
\draw[] (4) to node[w,above left] {24} (6);
\draw[] (6) to node[w,right] {35} (3);
\draw[] (2) to node[w,below right] {35} (3);
\draw[] (4) to node[w,below left] {33} (3);
\end{tikzpicture}
\end{center}
\end{pastexam}

\begin{pastsoln}
這份考卷和成大電機所 107 學年度的「資料結構」是同一份（兩系所共用），完整詳解見歷屆試題~\ref{px:ncku-107-ee-ds-09}。

答案：Kruskal 和 Sollin 得到同一棵樹，邊是 $5$--$6$ (10)、$0$--$1$ (20)、$4$--$6$ (24)、$1$--$6$ (32)、$3$--$4$ (33)、
$2$--$3$ (35)，總成本 $154$。 Sollin 第 1 輪加入 $0$--$1$、$2$--$3$、$3$--$4$、$4$--$6$、$5$--$6$，第 2 輪加入 $1$--$6$。
\end{pastsoln}
